\section{本章纲要：并立不是宋朝的背景}

辽与西夏不是夹在宋史年表两侧的边疆名词，而是各自经营都城、道路、税役、军队、文字与宗教机构的国家。辽把草原、东北与燕云的多中心网络接合起来；西夏则以河套、河西绿洲、城寨和多语文书维持权力。宋的银绢、边市、堡寨与使节制度，只有放在这两个行动者的资源、风险和国内政治中，才显出其含义。

本章不把和平当作战争之间的空白。盟约要靠交付、侦察、贸易管制和边地行政每日维持，冲突又会改变道路、迁徙和地方负担。读者将比较三方怎样把军事与财政延伸到不同区域，并通过辽的碑刻、墓志与城址，西夏的黑水城文书，以及宋辽夏之间的官修史和交聘记录，辨认制度名称、地方实践和史料视角之间的距离。
